% !TEX program = xelatex
% SKED document (Socratic Knowledge Elicitation & Documentation), phase 5.
% Build: python3 <plugin>/scripts/build.py --dir .   (or: xelatex main; xelatex main)
\documentclass[profile=ieee,lang=uk]{skedreport}

\skedtitle{Специфікація вимог до програмного забезпечення: Cinema Booking System}
\skedsubtitle{Внутрішній інженерний стандарт для розробки за методологією TDD}
\skedauthor{Бідник В.}
\skedaffiliation{}
\skedorganization{}
\skeddoctype{Software Requirements Specification (SRS)}
\skedversion{1.0.0}
\skeddate{\today}
\skedcity{Київ}
\skedudc{004.415.2}
\skedkeywords{SRS, кінотеатр, бронювання квитків, TDD, BDD, логування, атомарність, RBAC, IEEE 830}
\skedplatform{Вебзастосунок (доступ через веббраузер)}

\begin{document}

\skedmaketitle

\begin{abstract}
Специфікація вимог до програмного забезпечення (SRS) Cinema Booking System є внутрішнім документом для автора проєкту та AI-агента (K6, K10). Система реалізується як клієнт-серверний вебзастосунок (K9) із безумовною серверною валідацією прав доступу та бізнес-правил (K40) і персистентним збереженням даних (K41).

Архітектура визначає дві ролі: Користувач та єдиний початковий Адміністратор (K1, K27, K28). Перегляд списку фільмів і сеансів є публічним і доступний без входу (K2, K22). Для створення та скасування бронювань обов'язкові реєстрація та вхід за email і паролем (K23); користувач керує виключно власними бронюваннями за внутрішнім ID (K24). Без реєстрації створювати бронювання неможливо (гостьове бронювання та стороння автентифікація через OAuth/SSO у версії 1 не передбачені; K25). Онлайн-оплата винесена за межі системи (K33, K34); придбання квитків здійснюється офлайн (K35).

Обсяг одного бронювання становить від 1 до 6 місць (K4). Створення бронювання є строго атомарною операцією без проміжного утримання (hold) місць (K5, K31, K32): конкурентні конфлікти вирішуються на користь першого успішного запиту (K29), а часткове створення при зайнятості хоча б одного місця заборонене (K30). Скасування бронювання можливе лише повністю до моменту початку сеансу (K36, K37), переводить місця у вільні (K38) та є ідемпотентним (K39). Адміністратор керує сеансами, залами й фільмами (K3, K26): зміна параметрів сеансу (час, зал, фільм) дозволена за відсутності активних бронювань і блокується виключно за їх наявності (K44, K45); водночас адміністратор має право явно скасувати сеанс із фіксацією обов'язкової причини, відображенням її в кабінеті користувача та логуванням (K46, K47, K48) без автоматичного перенесення у v1 (K49).

Розробка базується на обов'язковому циклі TDD (K7, K11, K12, K13) без встановлення мінімального порогу покриття (K14). Система веде персистентні операційні та тестові логи без запису паролів і токенів (K8, K16, K17, K18, K19, K20). Вибір мови, фреймворків і СКБД зафіксовано як відкладене проєктне рішення з вимогами підтримки TDD, ролей та атомарності (K15, K42, K43, Q5).
\end{abstract}

\skedtoc

\section{Вступ}\label{sec:vstup}

\subsection{Призначення та цільова аудиторія}\label{sec:vstup-meta}
Цей документ являє собою внутрішню специфікацію вимог до програмного забезпечення (Software Requirements Specification, SRS) для інформаційної системи автоматизації бронювання квитків у кінотеатрі (Cinema Booking System)~\sked{K6}. Специфікація орієнтована на формат стандарту IEEE 830 (із врахуванням рекомендацій ISO/IEC/IEEE 29148) та призначена для фіксації погоджених вимог і обмежень проєкту між автором та автономним AI-агентом розробки Google Antigravity~\sked{K6,K10}.

Документ є єдиним джерелом правди для подальшої практичної реалізації системи за методологією Test-Driven Development (TDD) та верифікації її коректності за допомогою автоматизованих тестів~\sked{K10}.

\subsection{Межі системи та виключення версії 1}\label{sec:vstup-mezhi}
Система призначена для автоматизації процесів вибору сеансів та резервування місць у залах кінотеатру для двох ключових ролей: Користувача та Адміністратора~\sked{K1}. 

Для забезпечення високої надійності та швидкого постачання першої робочої версії (v1) низка функціональних можливостей свідомо винесена за межі системи:
\begin{itemize}
    \item \textbf{Гостьове бронювання та стороння автентифікація:} створення та скасування бронювань вимагають обов'язкової реєстрації у системі; гостьове бронювання та інтеграція із зовнішніми провайдерами автентифікації (OAuth, SSO, Google тощо) у першій версії не передбачені~\sked{K25}.
    \item \textbf{Оплата та фінансові транзакції:} система реалізує виключно безоплатне резервування місць; онлайн-оплата, інтеграція з платіжними сервісами (шлюзами), відстеження статусів оплати та процедури повернення коштів винесені за межі інформаційної системи~\sked{K33}.
    \item \textbf{Вартість квитків:} ціни квитків у першій версії не зберігаються, не обробляються базою даних та не відображаються у вебінтерфейсі користувача чи адміністратора~\sked{K34}.
    \item \textbf{Оплата квитків офлайн:} фактичне придбання квитків та їх оплата здійснюються користувачами безпосередньо в касах або терміналах кінотеатру поза межами цієї програми~\sked{K35}.
\end{itemize}

\subsection{Глосарій, терміни та скорочення}\label{sec:vstup-hlosarii}
У цій специфікації використовуються терміни та скорочення, наведені в \tabref{tab:glossary}~\sked{K6,K10}.

\begin{table}[htbp]
\caption{Глосарій ключових термінів та скорочень}\label{tab:glossary}
\centering
\begin{tabularx}{\textwidth}{@{}l l L@{}}
\toprule
\textbf{Термін/Скорочення} & \textbf{Повна назва} & \textbf{Визначення в контексті проєкту} \\
\midrule
SRS & Software Requirements Specification & Специфікація вимог до ПЗ відповідно до формату IEEE 830. \\
TDD & Test-Driven Development & Інженерна методологія розробки через попереднє написання автоматизованих тестів. \\
BDD & Behavior-Driven Development & Опис поведінки системи мовою бізнес-сценаріїв (Given-When-Then / Gherkin). \\
RBAC & Role-Based Access Control & Модель розмежування прав доступу на основі ролей облікового запису (Користувач, Адміністратор). \\
Фільм & Movie / Film & Кінокартина або художній твір, що має назву, тривалість, опис і демонструється в кінотеатрі. \\
Зал & Cinema Hall & Приміщення кінотеатру з визначеною назвою, місткістю та схемою розташування місць. \\
Сеанс & Showtime & Запланований показ фільму в залі з точним часом початку ($T_{сеансу}$). \\
Місце & Seat & Окрема одиниця розміщення глядача в залі (ряд і номер місця). \\
Бронювання & Booking / Reservation & Фіксація прав зареєстрованого користувача на перелік місць (від 1 до 6 включно) на обраний сеанс. \\
Атомарність & Atomicity & Властивість транзакції: створення бронювання виконується неподільно («все або нічого»). \\
Hold / Таймаут & Temporary Lock & Механізм тимчасового утримання місць на етапі вибору; у версії 1 не застосовується. \\
\bottomrule
\end{tabularx}
\end{table}

\subsection{Нормативні посилання та огляд структури}\label{sec:vstup-struktura}
Специфікація спирається на базові засади стандарту IEEE 830-1998 «IEEE Recommended Practice for Software Requirements Specifications» та проєктні домовленості, зафіксовані в репозиторії~\sked{K6}. 

Документ структуровано наступним чином: Розділ~\ref{sec:opis} наводить загальний архітектурний опис та обмеження; Розділ~\ref{sec:interfeisy} визначає контракт зовнішніх інтерфейсів на рівні операцій; Розділ~\ref{sec:funktsional} специфікує функціональні вимоги (REQ-F-xxx); Розділ~\ref{sec:tdd} встановлює обов'язковий регламент TDD і BDD-сценарії; Розділ~\ref{sec:loguvannia} регламентує правила аудиту та логування; Розділ~\ref{sec:nefunktsional} охоплює вимоги до якості та атомарності; Розділ~\ref{sec:vidkladeni} фіксує відкладені рішення щодо стеку; Розділ~\ref{sec:vysnovky} підсумовує готовність до розробки, а Додаток~\ref{sec:dodatok-matrytsia} містить матрицю простежуваності.

\section{Загальний опис системи}\label{sec:opis}

\subsection{Архітектура системи та перспектива продукту}\label{sec:opis-arkhitektura}
Система проєктується як сучасний вебзастосунок із класичною клієнт-серверною архітектурою~\sked{K9}. Клієнтська частина функціонує у веббраузері користувача та забезпечує відображення афіші, розкладу сеансів, схеми залу та інтерфейсу особистого кабінету~\sked{K9}. Серверна частина надає програмний інтерфейс (API), виконує бізнес-логіку та забезпечує транзакційну обробку операцій~\sked{K9}.

Ключовим архітектурним принципом є безумовна авторитетність серверної частини: перевірка прав доступу, коректності вхідних параметрів та бізнес-правил бронювання здійснюється на сервері незалежно від валідацій на стороні вебклієнта~\sked{K40}. Усі дані користувачів, фільмів, кінозалів, розкладу сеансів та створених бронювань зберігаються персистентно й гарантовано відновлюються після завершення роботи або перезапуску серверних процесів~\sked{K41}.

\begin{figure}[htbp]
\centering
\makebox[\textwidth][c]{%
\begin{tikzpicture}[node distance=10mm and 32mm]
  \node[sked box] (client) {Вебклієнт\\(Браузер / SPA)};
  \node[sked box, right=of client] (server) {Серверний застосунок\\(API \& Бізнес-правила)};
  \node[sked box, right=of server] (db) {Персистентне сховище\\(База даних / Сховище)};
  \draw[sked arrow] ([yshift=2.5mm]client.east) -- node[sked label, above=1pt] {Запити API} ([yshift=2.5mm]server.west);
  \draw[sked arrow] ([yshift=-2.5mm]server.west) -- node[sked label, below=1pt] {Відповіді \& Статуси} ([yshift=-2.5mm]client.east);
  \draw[sked arrow] ([yshift=2.5mm]server.east) -- node[sked label, above=1pt] {Атомарні транзакції} ([yshift=2.5mm]db.west);
  \draw[sked arrow] ([yshift=-2.5mm]db.west) -- node[sked label, below=1pt] {Персистентні дані} ([yshift=-2.5mm]server.east);
\end{tikzpicture}%
}
\caption{Контекстна клієнт-серверна архітектура системи}\label{fig:architecture}
\end{figure}

\subsection{Класи користувачів та характеристики ролей}\label{sec:opis-korystuvachi}
Система чітко розмежовує два класи користувачів на основі моделі рольового доступу (RBAC): Користувач (зареєстрований клієнт кінотеатру) та Адміністратор~\sked{K1}. 

\begin{itemize}
    \item \textbf{Користувач (клієнт):} зареєстрований відвідувач, що автентифікувався за email та паролем і має унікальний внутрішній числовий/ідентифікаційний ID. Має право переглядати репертуар, створювати бронювання на своє ім'я та переглядати чи скасовувати виключно власні бронювання~\sked{K1}. Звичайний користувач суворо позбавлений доступу до адміністративних операцій і не має технічної можливості призначити собі роль адміністратора~\sked{K27}.
    \item \textbf{Адміністратор:} співробітник кінотеатру, який працює в тій самій системі через спеціальний відокремлений розділ керування~\sked{K26}. Відповідає за наповнення довідників фільмів, залів і сеансів, а також за обробку позаштатних ситуацій (скасування сеансів)~\sked{K26}. Для першої версії передбачено створення рівно одного адміністративного облікового запису під час початкової конфігурації системи; підтримка керування кількома адміністраторами винесена за межі v1~\sked{K28}.
\end{itemize}

У \tabref{tab:rbac-matrix} наведено матрицю доступу до операцій системи для кожної ролі.

\begin{table}[htbp]
\caption{Матриця прав доступу ролей (RBAC)}\label{tab:rbac-matrix}
\centering
\begin{tabularx}{\textwidth}{@{}X >{\centering\arraybackslash}p{2.8cm} >{\centering\arraybackslash}p{2.6cm} >{\centering\arraybackslash}p{2.8cm}@{}}
\toprule
\textbf{Операція / Ресурс} & \textbf{Анонімний відвідувач} & \textbf{Користувач (Клієнт)} & \textbf{Адміністратор} \\
\midrule
Перегляд фільмів та сеансів & Дозволено & Дозволено & Дозволено \\
Перегляд схеми залу & Дозволено & Дозволено & Дозволено \\
Реєстрація та вхід (email+пароль) & Дозволено & Вже виконано & Вже виконано \\
Створення бронювання (1--6 місць) & Заборонено & Дозволено & Заборонено (у v1) \\
Перегляд власних бронювань & Заборонено & Дозволено & Дозволено \\
Скасування власних бронювань & Заборонено & Дозволено & Не застосовується \\
Керування фільмами, залами, сеансами & Заборонено & Заборонено & Дозволено \\
Явне скасування сеансу з причиною & Заборонено & Заборонено & Дозволено \\
Перегляд операційних логів & Заборонено & Заборонено & Дозволено \\
\bottomrule
\end{tabularx}
\end{table}

\subsection{Середовище функціонування та припущення}\label{sec:opis-seredovyshche}
Клієнтська частина повинна стабільно функціонувати в усіх сучасних веббраузерах (Google Chrome, Mozilla Firefox, Safari, Microsoft Edge). Серверна частина розгортається в середовищі, що підтримує роботу API та зв'язок зі сховищем даних~\sked{K9}.

Ключовим бізнес-припущенням є те, що бронювання слугує попереднім резервуванням місць, а безпосередній розрахунок грошима за квитки та отримання фіскальних чеків відбуваються в касі кінотеатру перед початком сеансу поза межами автоматизованої системи~\sked{K35}.

\subsection{Загальні обмеження проєктування}\label{sec:opis-obmezhennia}
При проєктуванні системи встановлюються такі засадничі обмеження:
\begin{itemize}
    \item \textbf{Ліміт бронювання:} одне бронювання повинно містити від 1 до 6 місць включно (максимально допустима кількість --- 6 місць)~\sked{K4}.
    \item \textbf{Відсутність утримання місць:} вибір місць на схемі залу у вебінтерфейсі не створює тимчасового блокування (hold) і не гарантує доступності місць до завершення атомарної операції запису~\sked{K31}.
    \item \textbf{Відмова від таймаутів у v1:} у версії 1 механізм блокування місць за таймаутом не реалізується; натомість застосовується перевірка конфліктів у момент оформлення~\sked{K32}.
\end{itemize}

\section{Вимоги до зовнішніх інтерфейсів}\label{sec:interfeisy}

Відповідно до погоджених проєктних рішень, зовнішні інтерфейси специфікуються на рівні необхідних бізнес-операцій та поведінки системи, без передчасної жорсткої фіксації конкретних URI REST-маршрутів чи JSON-схем, вибір яких здійснюється на етапі технічної реалізації~\sked{K40}.

\subsection{Операції вебінтерфейсу користувача}\label{sec:interfeisy-korystuvach}
Вебінтерфейс користувача забезпечує виконання наступних ключових операцій:
\begin{itemize}
    \item \textbf{Загальнодоступний перегляд (без входу):} користувач має змогу переглядати список актуальних фільмів, розклад сеансів на обрану дату та інтерактивну схему залу із позначенням зайнятих і вільних місць без необхідності автентифікації~\sked{K2,K22}.
    \item \textbf{Автентифікація та реєстрація:} форми введення облікових даних (email і пароль) для створення облікового запису та входу в систему.
    \item \textbf{Формування бронювання:} інтерактивний вибір від 1 до 6 вільних місць на схемі та ініціація дії «Забронювати». Інтерфейс повинен інформувати про результат операції (успіх або відхилення через конфлікт).
    \item \textbf{Особистий кабінет:} перегляд переліку власних бронювань із зазначенням фільму, сеансу, місць та поточного статусу («Активне», «Скасоване»). Якщо сеанс було скасовано адміністратором, інтерфейс обов'язково відображає зафіксовану причину скасування~\sked{K48}.
    \item \textbf{Скасування бронювання:} кнопка скасування активного бронювання з підтвердженням дії (доступна суто до часу початку сеансу).
\end{itemize}

\subsection{Операції вебінтерфейсу адміністратора}\label{sec:interfeisy-admin}
Адміністративний інтерфейс є окремим розділом системи, захищеним перевіркою ролі, та надає такі операційні можливості~\sked{K26}:
\begin{itemize}
    \item \textbf{Керування фільмами:} перегляд, додавання, редагування описів та деактивація фільмів.
    \item \textbf{Керування залами:} внесення інформації про кінозали та їх параметри.
    \item \textbf{Планування розкладу:} створення сеансів із прив'язкою до фільму, залу, дати та часу початку.
    \item \textbf{Модифікація сеансів:} редагування параметрів сеансів за відсутності бронювань. За наявності активних бронювань інтерфейс блокує зміну полів та надає окрему дію --- «Скасувати сеанс» із обов'язковим текстовим полем для введення причини.
\end{itemize}

\subsection{Сервісні операції та контракти поведінки бекенду}\label{sec:interfeisy-backend}
Серверна частина надає набір API-операцій, поведінка яких регулюється серверними контрактами (\tabref{tab:backend-operations})~\sked{K40}.

\begin{table}[htbp]
\caption{Логічні операції та контракти поведінки бекенду}\label{tab:backend-operations}
\centering
\begin{tabularx}{\textwidth}{@{}l l L@{}}
\toprule
\textbf{Логічна операція} & \textbf{Суб'єкт} & \textbf{Контракт поведінки та умови верифікації} \\
\midrule
Отримати розклад & Будь-хто & Повертає список доступних фільмів та сеансів без фільтрації за користувачем. \\
Отримати схему залу & Будь-хто & Повертає актуальний стан зайнятості місць на момент виконання запиту. \\
Реєстрація / Вхід & Анонім & Валідує формат email, перевіряє унікальність, видає сесійний ідентифікатор/токен. \\
Створити бронювання & Користувач & Атомарно перевіряє доступність усіх обраних місць (1--6) і фіксує бронювання. \\
Скасувати бронювання & Користувач & Перевіряє власника, статус та час сеансу; звільняє місця (ідемпотентна дія). \\
Створити сеанс & Адміністратор & Валідує відсутність накладання сеансів у залі, зберігає сеанс. \\
Скасувати сеанс & Адміністратор & Вимагає обов'язкової причини; каскадно переводить усі бронювання в «Скасовано». \\
\bottomrule
\end{tabularx}
\end{table}

\subsection{Загальні протокольні вимоги взаємодії}\label{sec:interfeisy-protokoly}
Взаємодія між клієнтом і сервером здійснюється за стандартними вебпротоколами (HTTP/HTTPS) із передачею структурованих повідомлень~\sked{K9}. Сервер повинен повертати чіткі статус-коди результату операції (успіх, помилка авторизації, бізнес-конфлікт зайнятих місць, помилка валідації) та зрозумілі повідомлення про причини відмови~\sked{K40}.

\section{Функціональні вимоги}\label{sec:funktsional}

Усі функціональні вимоги системи мають стабільні унікальні ідентифікатори формату \texttt{REQ-F-xxx}, супроводжуються посиланнями на консенсусний простір знань $K^*$ та містять чіткі умови верифікації.

\subsection{Каталог фільмів і розклад сеансів}\label{sec:funktsional-katalog}

\begin{itemize}[leftmargin=*,labelindent=0pt]
    \item \textbf{REQ-F-001 (Загальнодоступний перегляд розкладу):} Система повинна надавати можливість перегляду актуального списку фільмів та сеансів будь-якому відвідувачеві без проходження автентифікації.~\sked{K2,K22}
    \textit{Умова верифікації:} Запит списку фільмів і сеансів без сесійного токена повертає успішний статус та коректні дані розкладу.
    
    \item \textbf{REQ-F-002 (Відображення схеми залу):} Система повинна відображати схему кінозалу обраного сеансу із актуальним статусом кожного місця («вільне» або «зайняте»). Вибір місця у вебінтерфейсі не резервує його і не блокує для інших користувачів.~\sked{K2,K31}
    \textit{Умова верифікації:} Вибір місця одним користувачем у браузері не змінює статус місця для інших користувачів до моменту підтвердження бронювання.
\end{itemize}

\subsection{Автентифікація, реєстрація та управління доступом}\label{sec:funktsional-auth}

\begin{itemize}[leftmargin=*,labelindent=0pt]
    \item \textbf{REQ-F-003 (Реєстрація та вхід за email):} Для створення та скасування бронювань користувач повинен зареєструватися й увійти в систему за допомогою дійсного email та пароля.~\sked{K23}
    \textit{Умова верифікації:} Спроба створення або скасування бронювання без автентифікації повертає статус відмови в авторизації.
    
    \item \textbf{REQ-F-004 (Внутрішній ID користувача):} Кожен зареєстрований обліковий запис отримує унікальний внутрішній ID, який пов'язується з усіма його транзакціями та логами.~\sked{K24}
    \textit{Умова верифікації:} Внутрішній ID користувача генерується автоматично, є незмінним та унікальним у сховищі даних.
    
    \item \textbf{REQ-F-005 (Ізоляція доступу до бронювань):} Користувач має право переглядати та скасовувати виключно власні бронювання.~\sked{K24}
    \textit{Умова верифікації:} Спроба перегляду або скасування бронювання користувача А за ідентифікатором користувача Б повертає статус заборони доступу (Forbidden).
    
    \item \textbf{REQ-F-006 (Відокремлення прав адміністратора):} Адміністратор працює в системі через окремий розділ керування довідниками. Звичайний користувач суворо позбавлений доступу до адміністративного розділу.~\sked{K26}
    \textit{Умова верифікації:} Запит користувача з роллю «Користувач» на адміністративний маршрут повертає статус заборони доступу (Forbidden).
\end{itemize}

\subsection{Створення бронювання та атомарне запобігання конфліктам}\label{sec:funktsional-bronuvannia}

\begin{figure}[htbp]
\centering
\begin{tikzpicture}[node distance=10mm and 12mm]
  \node[sked box] (free) {Вільне місце\\в залі};
  \node[sked box, right=of free] (booked) {Статус\\«Активне»};
  \node[sked box, right=of booked] (cancelled) {Статус\\«Скасоване»};
  \draw[sked arrow] (free) -- node[sked label, above=2pt] {Атомарне бронювання} (booked);
  \draw[sked arrow] (booked) -- node[sked label, above=2pt] {Скасування (до сеансу)} (cancelled);
  \draw[sked back] ([yshift=-4mm]cancelled.south) to[out=-135,in=-45] 
        node[sked label, below=2pt] {Місця знову вільні} ([yshift=-4mm]free.south);
\end{tikzpicture}
\caption{Життєвий цикл стану бронювання та місця}\label{fig:booking-lifecycle}
\end{figure}

\begin{itemize}[leftmargin=*,labelindent=0pt]
    \item \textbf{REQ-F-007 (Діапазон кількості місць):} Одне бронювання повинно містити від 1 до 6 місць включно.~\sked{K4}
    \textit{Умова верифікації:} Запит на створення бронювання з 0 або понад 6 місцями відхиляється з повідомленням про порушення обмеження діапазону.
    
    \item \textbf{REQ-F-008 (Атомарне створення без утримання):} Перевірка доступності всіх обраних місць та створення бронювання повинні виконуватися як єдина неподільна транзакція в момент натискання кнопки «Забронювати». Тимчасове утримання (hold) не застосовується.~\sked{K5}
    \textit{Умова верифікації:} У базі даних створення запису бронювання та маркування місць виконуються в одній транзакції без проміжних станів.
    
    \item \textbf{REQ-F-009 (Вирішення конкурентних конфліктів):} При одночасній спробі кількох користувачів забронювати одне й те саме місце на той самий сеанс успішним може бути лише одне бронювання. Усі інші запити відхиляються з повідомленням про зайнятість місця.~\sked{K29}
    \textit{Умова верифікації:} При імітації 10 паралельних запитів на одне місце успішним стає рівно один, а 9 запитів отримують відмову зі статусом конфлікту через зайнятість місця.
    
    \item \textbf{REQ-F-010 (Принцип «все або нічого»):} Якщо бронювання містить кілька місць і хоча б одне з них уже зайняте іншим користувачем, усе бронювання відхиляється повністю. Часткове бронювання суворо заборонене.~\sked{K30}
    \textit{Умова верифікації:} Запит на 3 місця, де 2 вільні й 1 зайняте, відхиляється без створення часткового бронювання; усі вільні місця залишаються доступними.
\end{itemize}

\subsection{Скасування бронювання користувачем}\label{sec:funktsional-skasuvannia-user}

\begin{itemize}[leftmargin=*,labelindent=0pt]
    \item \textbf{REQ-F-011 (Повне скасування бронювання):} Користувач може скасувати лише власне бронювання і виключно в повному обсязі (усі місця в ньому одночасно). Часткове скасування окремих місць заборонене.~\sked{K36}
    \textit{Умова верифікації:} Запит на скасування переводить усе бронювання у статус скасованого, вибіркові місця не скасовуються.
    
    \item \textbf{REQ-F-012 (Часовий дедлайн скасування):} Скасування бронювання дозволене виключно до настання часу початку сеансу ($T_{поточний} < T_{сеансу}$). У момент початку сеансу та після нього операція скасування блокується системою.~\sked{K37}
    \textit{Умова верифікації:} Запит на скасування, виконаний на 1 секунду пізніше за час початку сеансу, відхиляється з кодом помилки.
    
    \item \textbf{REQ-F-013 (Звільнення місць після скасування):} Після успішного скасування бронювання отримує статус «Скасоване», а всі зарезервовані за ним місця негайно стають доступними для вибору іншими користувачами.~\sked{K38}
    \textit{Умова верифікації:} Після успішного скасування повторний запит схеми залу відображає ці місця як вільні для нового бронювання.
    
    \item \textbf{REQ-F-014 (Ідемпотентність повторного скасування):} Повторний запит на скасування вже скасованого бронювання не змінює стан системи, не генерує побічних ефектів та є безпечним.~\sked{K39}
    \textit{Умова верифікації:} Повторний виклик скасування повертає інформацію про скасований статус без помилки та без повторного звільнення місць.
\end{itemize}

\subsection{Керування та скасування сеансів адміністратором}\label{sec:funktsional-admin}

\begin{itemize}[leftmargin=*,labelindent=0pt]
    \item \textbf{REQ-F-015 (Керування довідниками):} Адміністратор здійснює базові операції внесення та редагування інформації про фільми, кінозали та розклад сеансів.~\sked{K3}
    \textit{Умова верифікації:} Адміністратор може успішно створити новий сеанс для наявного фільму та кінозалу.
    
    \item \textbf{REQ-F-016 (Вільне редагування сеансів без бронювань):} Якщо на сеанс немає жодного активного бронювання, адміністратор має змогу вільно змінювати час початку, зал або фільм.~\sked{K44}
    \textit{Умова верифікації:} Редагування часу сеансу за відсутності бронювань завершується успішно з оновленням розкладу.
    
    \item \textbf{REQ-F-017 (Блокування зміни сеансів з бронюваннями):} Якщо на сеанс створено хоча б одне активне бронювання, пряме редагування параметрів сеансу (час, зал, фільм) повністю блокується системою.~\sked{K45}
    \textit{Умова верифікації:} Спроба оновлення часу або залу сеансу з активними бронюваннями повертає статус конфлікту через наявність активних бронювань.
    
    \item \textbf{REQ-F-018 (Явне скасування сеансу з причиною):} За наявності активних бронювань адміністратор має право явно скасувати сеанс за умови обов'язкового зазначення причини скасування та явного підтвердження дії.~\sked{K46}
    \textit{Умова верифікації:} Запит на скасування сеансу без заповненого поля причини відхиляється валідатором.
    
    \item \textbf{REQ-F-019 (Каскадне скасування бронювань):} При скасуванні сеансу сам сеанс та всі пов'язані з ним активні бронювання отримують статус «Скасовано» й зберігаються в історії для аудиту.~\sked{K47}
    \textit{Умова верифікації:} Після скасування сеансу всі бронювання цього сеансу в базі даних автоматично мають статус «Скасовано».
    
    \item \textbf{REQ-F-020 (Відображення причини в кабінеті користувача):} Користувач повинен бачити факт скасування свого бронювання адміністратором та текстову причину скасування сеансу у власному особистому кабінеті.~\sked{K48}
    \textit{Умова верифікації:} У відповіді API списку бронювань користувача для скасованого сеансу повертається поле з причиною від адміністратора.
\end{itemize}

\section{Інженерний процес TDD, критерії прийняття та BDD-сценарії}\label{sec:tdd}

Цей розділ фіксує обов'язковий інженерний регламент розробки програмного забезпечення за методологією Test-Driven Development (TDD) та формальні верифікаційні критерії прийняття.

\subsection{Обов'язковий життєвий цикл розробки TDD}\label{sec:tdd-process}

\begin{figure}[htbp]
\centering
\begin{tikzpicture}[node distance=8mm and 10mm]
  \node[sked box] (t1) {1. Написання\\тесту};
  \node[sked box, right=of t1] (t2) {2. Перевірка\\падіння (Red)};
  \node[sked box, right=of t2] (t3) {3. Реалізація\\коду (Green)};
  \node[sked box, below=of t3] (t4) {4. Рефакторинг\\коду};
  \node[sked box, left=of t4] (t5) {5. Повторний\\запуск тестів};
  
  \draw[sked arrow] (t1) -- (t2);
  \draw[sked arrow] (t2) -- (t3);
  \draw[sked arrow] (t3) -- (t4);
  \draw[sked arrow] (t4) -- (t5);
  \draw[sked back] (t5) to[out=180,in=180,looseness=1.5] node[sked label, left=2pt] {Наступна функція} (t1);
\end{tikzpicture}
\caption{Обов'язковий 5-кроковий цикл розробки за TDD}\label{fig:tdd-cycle}
\end{figure}

\begin{itemize}[leftmargin=*,labelindent=0pt]
    \item \textbf{REQ-TDD-001 (Обов'язковість процесу TDD):} Розробка будь-якої нової функціональності, бізнес-правил та виправлення дефектів повинні здійснюватися суворо за методологією TDD.~\sked{K7,K11}
    \textit{Умова верифікації:} Будь-який комміт або пул-реквест у репозиторій повинен містити відповідні автоматизовані тести, створені суворо перед написанням коду реалізації.

    \item \textbf{REQ-TDD-002 (5-кроковий цикл розробки):} Кожен етап розробки виконується за чіткою послідовністю кроків: (1) створити автоматизований тест для очікуваної поведінки перед написанням функціонального коду; (2) переконатися, що тест падає саме з очікуваної причини (відсутність функції/неправильний результат); (3) написати мінімальний код для проходження тесту; (4) провести рефакторинг коду; (5) повторно запустити відповідний набір тестів і переконатися в їх успішності.~\sked{K12}
    \textit{Умова верифікації:} Історія змін та логи запусків тестів підтверджують проходження стадій «Red $\rightarrow$ Green $\rightarrow$ Refactor».
\end{itemize}

\subsection{Сфера тестування та перевірка граничних випадків}\label{sec:tdd-scope}

\begin{itemize}[leftmargin=*,labelindent=0pt]
    \item \textbf{REQ-TDD-003 (Покриття бізнес-правил та граничних випадків):} Автоматизовані тести повинні верифікувати всі погоджені в цій SRS бізнес-правила, обмеження ролей та суттєві граничні випадки системи бронювання.~\sked{K13}
    \textit{Обов'язкові граничні тест-кейси:}
    \begin{itemize}
        \item Бронювання рівно 1 місця (нижня межа) та рівно 6 місць (верхня межа);
        \item Спроба бронювання 0 та 7 місць (порушення меж діапазону);
        \item Конкурентна спроба забронювання одного й того самого місця кількома паралельними потоками;
        \item Запит на бронювання кількох місць, коли одне з них уже зайняте;
        \item Скасування за 1 секунду до часу початку сеансу (успіх) та в момент початку (відхилення);
        \item Повторний виклик скасування вже скасованого бронювання (перевірка ідемпотентності);
        \item Спроба модифікації сеансу з активними бронюваннями (блокування) та скасування сеансу без введення причини (блокування).
    \end{itemize}
    
    \item \textbf{REQ-TDD-004 (Відсутність жорсткого порогу покриття):} Конкретний фіксований відсоток покриття рядків коду тестами (coverage threshold) на поточному етапі не встановлюється; пріоритетом є якісна повнота перевірки бізнес-сценаріїв та граничних умов.~\sked{K14}
    
    \item \textbf{REQ-TDD-005 (Критерії до тестового інструментарію):} Фреймворк тестування обирається разом із технологічним стеком, проте обов'язково повинен підтримувати швидкий запуск юніт- та інтеграційних тестів, можливість емуляції конкурентних запитів та генерацію структурованих звітів запусків.~\sked{K15,K43}
\end{itemize}

\subsection{Критерії прийняття та BDD-сценарії поведінки системи}\label{sec:tdd-bdd}
Для ключових бізнес-правил формулюються обов'язкові сценарії мовою Gherkin (Given-When-Then), що слугують безпосередньою основою для написання тестів~\sked{K13}.

\subsubsection{Сценарій 1: Конкурентне бронювання одного місця}\label{sec:bdd-sc1}
\begin{skedcode}
Сценарій: Конкурентне бронювання одного й того самого місця двома користувачами
  Дано: Користувач_1 та Користувач_2 автентифіковані у системі
    Та місце "Ряд 3, Місце 5" на сеанс ID=101 є вільним
  Коли: Користувач_1 надсилає запит на бронювання місця "Ряд 3, Місце 5"
    Та одночасно Користувач_2 надсилає запит на бронювання того ж місця
  Тоді: Рівно одне бронювання створюється зі статусом "Активне"
    Та інший користувач отримує помилку про зайнятість місця
    Та місце залишається зайнятим і не бронюється повторно
\end{skedcode}
\noindent\textit{Пов'язані вимоги:} REQ-F-008, REQ-F-009, REQ-NF-001~\sked{K29}.

\subsubsection{Сценарій 2: Відхилення бронювання при зайнятості хоча б одного місця}\label{sec:bdd-sc2}
\begin{skedcode}
Сценарій: Бронювання трьох місць, коли одне з них зайняте (все або нічого)
  Дано: Користувач автентифікований у системі
    Та на сеанс ID=101 місця 1 і 2 вільні, а місце 3 уже заброньоване іншим клієнтом
  Коли: Користувач надсилає запит на бронювання групи місць [1, 2, 3]
  Тоді: Запит відхиляється повністю
    Та жодного запису бронювання для користувача не створюється
    Та місця 1 і 2 залишаються вільними для вибору
\end{skedcode}
\noindent\textit{Пов'язані вимоги:} REQ-F-010, REQ-NF-001~\sked{K30}.

\subsubsection{Сценарій 3: Успішне скасування бронювання до початку сеансу}\label{sec:bdd-sc3}
\begin{skedcode}
Сценарій: Успішне скасування бронювання до початку сеансу
  Дано: Користувач автентифікований у системі
    Та користувач має початкове активне бронювання ID=501 на сеанс із часом початку 19:00
    Та поточний системний час становить 18:45 (менший за час початку сеансу)
  Коли: Користувач ініціює скасування власного бронювання ID=501
  Тоді: Бронювання переходить у статус "Скасоване"
    Та всі зарезервовані за ним місця негайно стають вільними
\end{skedcode}
\noindent\textit{Пов'язані вимоги:} REQ-F-011, REQ-F-012, REQ-F-013~\sked{K36,K37,K38}.

\subsubsection{Сценарій 4: Блокування скасування бронювання точно в момент початку сеансу}\label{sec:bdd-sc4}
\begin{skedcode}
Сценарій: Блокування скасування точно в момент настання часу початку сеансу
  Дано: Користувач автентифікований у системі
    Та користувач має початкове активне бронювання ID=502 на сеанс із часом початку 19:00:00
    Та поточний системний час становить рівно 19:00:00 (настав час початку сеансу)
  Коли: Користувач надсилає запит на скасування власного бронювання ID=502
  Тоді: Система блокує операцію з повідомленням про початок сеансу
    Та статус бронювання ID=502 залишається "Активне"
    Та зарезервовані місця не звільняються
\end{skedcode}
\noindent\textit{Пов'язані вимоги:} REQ-F-012~\sked{K37}.

\subsubsection{Сценарій 5: Блокування скасування бронювання після початку сеансу}\label{sec:bdd-sc5}
\begin{skedcode}
Сценарій: Блокування спроби скасування бронювання після початку сеансу
  Дано: Користувач автентифікований у системі
    Та користувач має початкове активне бронювання ID=503 на сеанс із часом початку 19:00
    Та поточний системний час становить 19:15 (сеанс уже триває)
  Коли: Користувач намагається скасувати власне бронювання ID=503
  Тоді: Система блокує операцію з повідомленням про те, що сеанс уже триває
    Та статус бронювання ID=503 залишається "Активне"
    Та зарезервовані місця не звільняються
\end{skedcode}
\noindent\textit{Пов'язані вимоги:} REQ-F-012~\sked{K37}.

\subsubsection{Сценарій 6: Блокування зміни та явне скасування сеансу адміністратором}\label{sec:bdd-sc6}
\begin{skedcode}
Сценарій: Обробка сеансу з активними бронюваннями адміністратором
  Дано: Адміністратор увійшов до адміністративного розділу
    Та на сеанс ID=202 існують початкові активні бронювання користувачів
  Коли: Адміністратор намагається змінити час або зал сеансу ID=202
  Тоді: Система блокує зміну параметрів сеансу через наявність активних бронювань
  Коли: Адміністратор обирає "Скасувати сеанс", вводить причину "Технічна несправність проєктора" і підтверджує
  Тоді: Сеанс ID=202 отримує статус "Скасовано"
    Та всі активні бронювання на сеанс ID=202 переходять у статус "Скасовано"
    Та в кабінеті кожного клієнта відображається причина "Технічна несправність проєктора"
\end{skedcode}
\noindent\textit{Пов'язані вимоги:} REQ-F-017, REQ-F-018, REQ-F-019, REQ-F-020~\sked{K45,K46}.

\section{Вимоги до логування та аудиту}\label{sec:loguvannia}

Система веде два незалежні рівні журналювання: (1) операційні логи бізнес-подій у процесі функціонування застосунку та (2) протоколи виконання автоматизованих тестів у процесі TDD~\sked{K8}. Усі вимоги мають ідентифікатори \texttt{REQ-LOG-xxx}.

\subsection{Операційне логування подій}\label{sec:loguvannia-podii}

\begin{itemize}[leftmargin=*,labelindent=0pt]
    \item \textbf{REQ-LOG-001 (Обов'язкові типи подій для логування):} Система повинна автоматично фіксувати у журналі такі типи подій:
    \begin{enumerate}
        \item Створення нового бронювання користувачем;
        \item Скасування бронювання користувачем;
        \item Модифікація довідників адміністратором (створення, оновлення або скасування фільмів, кінозалів і сеансів);
        \item Відмови в операціях через конфлікти зайнятих місць або порушення правил бронювання;
        \item Непередбачені системні збої та внутрішні помилки застосунку.~\sked{K16}
    \end{enumerate}
    \textit{Умова верифікації:} Кожна з перелічених дій ініціює запис відповідного рядка у журнал аудиту.

    \item \textbf{REQ-LOG-002 (Структура та обов'язковий склад полів запису):} Кожен запис операційного журналу повинен бути структурованим (наприклад, JSON-об'єкт) та містити обов'язкові поля, наведені у \tabref{tab:log-fields}~\sked{K17}.
\end{itemize}

\begin{table}[htbp]
\caption{Обов'язкові атрибути операційного лог-запису}\label{tab:log-fields}
\centering
\begin{tabularx}{\textwidth}{@{}l l L@{}}
\toprule
\textbf{Поле} & \textbf{Тип даних} & \textbf{Опис та призначення} \\
\midrule
\texttt{timestamp} & ISO-8601 UTC & Точна дата і час фіксації події з мілісекундами. \\
\texttt{event\_type} & String (Enum) & Тип події (\texttt{BOOKING\_CREATED}, \texttt{SESSION\_CANCELLED} тощо). \\
\texttt{status} & String & Результат операції (\texttt{SUCCESS}, \texttt{REJECTED}, \texttt{ERROR}). \\
\texttt{user\_id} & Integer / UUID & Внутрішній ID користувача або адміністратора (якщо ідентифікований). \\
\texttt{entity\_ids} & Object & Ідентифікатори об'єктів (ID бронювання, ID сеансу, номери місць). \\
\texttt{reason} & String & Детальна текстова причина відмови для неуспішних операцій. \\
\bottomrule
\end{tabularx}
\end{table}

\subsection{Захист конфіденційності та правила фільтрації логів}\label{sec:loguvannia-bezpeka}

\begin{itemize}[leftmargin=*,labelindent=0pt]
    \item \textbf{REQ-LOG-003 (Заборона збереження конфіденційних даних):} Система у версії 1 не збирає і не обробляє платіжні реквізити. До журналів логів суворо заборонено записувати паролі у відкритому чи хешованому вигляді, токени доступу (сесійні токени/JWT) та будь-які інші конфіденційні автентифікаційні дані~\sked{K18}.
    \textit{Умова верифікації:} Автоматизований тест-фільтр перевіряє вихідний потік логів під час реєстрації та входу; поля паролів замінюються маскою \texttt{[REDACTED]} або повністю вилучаються.
\end{itemize}

\subsection{Персистентність та зберігання логів для діагностики}\label{sec:loguvannia-zberezhennia}

\begin{itemize}[leftmargin=*,labelindent=0pt]
    \item \textbf{REQ-LOG-004 (Персистентність операційних логів):} Логи повинні записуватися на персистентний носій і зберігатися після штатного чи аварійного завершення роботи програми для можливості подальшого аналізу та діагностики помилок. Конкретний механізм збереження (структуровані файли, ротація, база даних) обирається під час проєктування технічної реалізації~\sked{K19}.
    \textit{Умова верифікації:} Після примусового перезавантаження серверного процесу журнал логів попередніх операцій зберігає цілісність і доступний для читання.
\end{itemize}

\subsection{Журналювання результатів автоматизованого тестування}\label{sec:loguvannia-testy}

\begin{itemize}[leftmargin=*,labelindent=0pt]
    \item \textbf{REQ-LOG-005 (Збереження протоколів запусків тестів TDD):} Результати кожного запуску автоматизованих тестів повинні зберігатися в окремих журналах (тест-логах). Журнал запусків фіксує: час початку й завершення тестування, загальну кількість тестів, повний перелік успішно пройдених і впалих тест-кейсів, а для кожного падіння --- детальний трасування стеку та діагностичне повідомлення про причину падіння~\sked{K20}.
    \textit{Умова верифікації:} Тестовий раннер автоматично генерує стандартизований звіт після кожного тестового прогону.
\end{itemize}

\section{Нефункціональні вимоги та системні обмеження}\label{sec:nefunktsional}

Усі нефункціональні вимоги мають ідентифікатори формату \texttt{REQ-NF-xxx} і зосереджені на якісних характеристиках надійності, безпеки та транзакційної цілісності даних.

\subsection{Транзакційна атомарність та надійність даних}\label{sec:nefunktsional-atomarnist}

\begin{itemize}[leftmargin=*,labelindent=0pt]
    \item \textbf{REQ-NF-001 (Транзакційна атомарність створення бронювання):} Система повинна гарантувати сувору атомарність (властивість ACID) для операції створення бронювання. Будь-який збій на проміжному етапі повинен повністю відкочувати транзакцію до попереднього узгодженого стану (принцип «все або нічого»)~\sked{K5,K30}.
    \textit{Умова верифікації:} Штучне переривання транзакції запису в базу даних не залишає частково заброньованих місць або бронювань без прив'язки до місць.

    \item \textbf{REQ-NF-002 (Персистентність та збереження цілісності):} Серверна частина повинна гарантувати персистентне збереження та цілісність усіх сутностей (користувачі, фільми, зали, сеанси, бронювання). Рестарт системи чи падіння окремих процесів не повинні призводити до втрати узгодженості або дублювання даних~\sked{K40,K41}.
    \textit{Умова верифікації:} Автоматичний тест перезавантаження сервера підтверджує 100\% збереження всіх зафіксованих перед зупинкою бронювань.
\end{itemize}

\subsection{Безпека та захист від несанкціонованого доступу}\label{sec:nefunktsional-bezpeka}

\begin{itemize}[leftmargin=*,labelindent=0pt]
    \item \textbf{REQ-NF-003 (Серверний контроль прав RBAC):} Усі операції зміни даних та доступу до приватних ресурсів повинні перевірятися на рівні бізнес-логіки сервера. Клієнтські перевірки в браузері розглядаються виключно як засіб підвищення зручності інтерфейсу, але не як засіб безпеки~\sked{K40}.
    \textit{Умова верифікації:} Запити в обхід вебклієнта (через прямі виклики API) суворо підпорядковуються серверним правилам авторизації.

    \item \textbf{REQ-NF-004 (Ізоляція даних між користувачами):} Система повинна унеможливлювати витік інформації про бронювання іншого клієнта чи маніпулювання чужими бронюваннями (захист від атак типу Insecure Direct Object References, IDOR)~\sked{K24}.
    \textit{Умова верифікації:} Зміна ID бронювання в тілі або параметрах запиту на ID чужого бронювання повертає помилку доступу.

    \item \textbf{REQ-NF-005 (Неможливість самопризначення ролі адміністратора):} Процедура реєстрації звичайного користувача повинна створювати обліковий запис виключно з роллю «Користувач». Спроби передачі адміністративних прапорців у запиті реєстрації повинні ігноруватися або відхилятися~\sked{K27}.
    \textit{Умова верифікації:} Перевірка створеного запису в базі даних після реєстрації із запитом, що містить поле \texttt{role: "admin"}, підтверджує присвоєння ролі звичайного користувача.
\end{itemize}

\subsection{Обробка конкурентних запитів та продуктивність}\label{sec:nefunktsional-konkurentnist}

\begin{itemize}[leftmargin=*,labelindent=0pt]
    \item \textbf{REQ-NF-006 (Коректність при конкурентному навантаженні):} Відповідно до узгоджених вимог, для першої версії не встановлюються штучні числові пороги часу відповіді чи кількості одночасних користувачів. Головною вимогою до продуктивності є якісна здатність системи надійно опрацьовувати одночасні запити до одного й того самого сеансу без взаємних блокувань (дедлоків), деградації цілісності бази даних та гарантувати, що одне місце ні за яких умов не буде заброньоване двома різними клієнтами~\sked{K29,K43}.
    \textit{Умова верифікації:} Стрес-тест конкурентного доступу (100 паралельних запитів на одні й ті самі місця) підтверджує відсутність подвійних бронювань та завершення транзакцій без дедлоків.
\end{itemize}

\section{Відкладені проєктні рішення та виключення версії 1}\label{sec:vidkladeni}

У цьому розділі зафіксовано свідомі обмеження поточної специфікації: інженерні рішення, вибір яких відкладено на етап проєктування реалізації, а також функції, які явно виключено з першої версії системи.

\subsection{Відкладені рішення щодо технологічного стеку та фреймворків}\label{sec:vidkladeni-stek}
Конкретна мова програмування, вебфреймворк для побудови API, система керування базами даних (СКБД) та бібліотеки для автоматизованого тестування у цій специфікації навмисно не фіксуються~\sked{K42,Q5}. Їх вибір здійснюється розробником на наступному етапі проєктування реалізації та погоджується перед початком написання коду~\sked{K15,K42,Q5}.

При цьому до обраного технологічного стеку встановлюються чіткі обов'язкові критерії (\tabref{tab:tech-criteria})~\sked{K43,Q5}.

\begin{table}[htbp]
\caption{Обов'язкові критерії до вибору технологій реалізації}\label{tab:tech-criteria}
\centering
\begin{tabularx}{\textwidth}{@{}l L@{}}
\toprule
\textbf{Вимога / Критерій} & \textbf{Технічне обґрунтування та значення для системи} \\
\midrule
Підтримка TDD & Наявність зрілих інструментів для юніт- та інтеграційного тестування з високою швидкістю виконання тест-сьютів. \\
Рольове розмежування & Ефективні вбудовані засоби автентифікації та перевірки прав доступу на рівні бекенду (RBAC). \\
Транзакційна атомарність & Підтримка транзакційного механізму в базі даних (ACID / атомарні операції) для гарантії захисту від подвійних бронювань. \\
Персистентність даних & Надійне збереження даних на диску з можливістю коректного відновлення після рестарту. \\
\bottomrule
\end{tabularx}
\end{table}

\subsection{Погоджені виключення та обмеження першої версії}\label{sec:vidkladeni-vykliuchennia}
Відповідно до рішень автора проєкту, у версії 1 не реалізуються наступні функціональні можливості (\tabref{tab:v1-exclusions})~\sked{K25,K28,K32,K33,K34,K49}.

\begin{table}[htbp]
\caption{Погоджені функціональні виключення першої версії}\label{tab:v1-exclusions}
\centering
\begin{tabularx}{\textwidth}{@{}l L L@{}}
\toprule
\textbf{Функціональна область} & \textbf{Виключення у версії 1} & \textbf{Погоджений статус v1} \\
\midrule
Автентифікація & Гостьове бронювання та сторонній вхід (Google, OAuth, SSO). & Лише реєстрація через email/пароль. \\
Оплата квитків & Онлайн-оплата, платіжні шлюзи, статуси оплат, повернення коштів. & Безоплатне резервування (оплата в касі). \\
Ціноутворення & Збереження, розрахунок та відображення цін квитків. & Ціни не відображаються у системі. \\
Бронювання & Тимчасове утримання місць (hold) із таймаутом на час вибору. & Атомарна фіксація в момент кліку. \\
Адміністрування & Керування кількома адміністраторами та правами адмінів. & Єдиний системний адміністратор. \\
Зміна сеансів & Автоматичне перенесення бронювань та email/SMS розсилки. & Ручне перебронювання користувачами. \\
\bottomrule
\end{tabularx}
\end{table}

\section{Висновки}\label{sec:vysnovky}

\subsection{Підсумок вимог та критерії готовності до реалізації}\label{sec:vysnovky-pidsumok}
Ця специфікація вимог до програмного забезпечення (SRS) фіксує узгоджений консенсусний простір знань $K^*$ для Cinema Booking System, сформований у ході ітеративного сократичного діалогу між автором проєкту та автономним AI-агентом Google Antigravity~\sked{K6,K10}.

Специфікація повністю покриває всі критичні аспекти функціонування системи:
\begin{itemize}
    \item Чітке розмежування ролей Користувача та єдиного системного Адміністратора;
    \item Загальнодоступний перегляд сеансів та захищене створення/скасування бронювань виключно для зареєстрованих клієнтів;
    \item Строге дотримання обмеження від 1 до 6 місць на одне бронювання;
    \item Надійну атомарну модель захисту від подвійних бронювань без застосування складних таймаутів (hold) у версії 1;
    \item Правило «все або нічого» при бронюванні групи місць;
    \item Однозначні правила скасування бронювань користувачем (лише повне скасування суто до часу початку сеансу);
    \item Правила модифікації розкладу адміністратором (блокування редагування сеансів з активними бронюваннями та можливість явного скасування із зазначенням обов'язкової причини);
    \item Обов'язковість дисципліни Test-Driven Development (TDD) як засадничого процесу розробки нової функціональності та усунення дефектів~\sked{K7};
    \item Регламент операційного та тестового логування із гарантією захисту персональних і конфіденційних даних.
\end{itemize}

Зафіксований набір вимог та BDD-сценаріїв у Додатку~\ref{sec:dodatok-matrytsia} є вичерпним і достатнім для переходу до етапу інженерного проєктування та безпосередньої реалізації коду системи за методологією TDD~\sked{K7,K10}.


\skedappendix
\section{Початкова матриця простежуваності вимог}\label{sec:dodatok-matrytsia}

Цей додаток містить початкову матрицю простежуваності (Traceability Matrix), яка забезпечує двосторонній зв'язок між зафіксованими пунктами погодженого консенсусу $K^*$, формальними ідентифікаторами вимог (\texttt{REQ-F}, \texttt{REQ-TDD}, \texttt{REQ-LOG}, \texttt{REQ-NF}), конкретними умовами їх верифікації та відповідними BDD-сценаріями перевірки поведінки (\tabref{tab:traceability}).

\begin{small}
\setlength{\tabcolsep}{3.5pt}
\begin{longtable}{@{}p{2.2cm} p{1.4cm} p{1.6cm} p{7.4cm} p{2.4cm}@{}}
\caption{Матриця простежуваності: Вимоги $\rightarrow$ База $K^*$ $\rightarrow$ Критерії $\rightarrow$ BDD-сценарії}\label{tab:traceability}\\
\toprule
\textbf{ID вимоги} & \textbf{Пункт $K^*$} & \textbf{Розділ SRS} & \textbf{Критерій прийняття} & \textbf{BDD-сценарій} \\
\midrule
\endfirsthead
\multicolumn{5}{l}{\small\itshape Продовження таблиці~\ref{tab:traceability}}\\
\toprule
\textbf{ID вимоги} & \textbf{Пункт $K^*$} & \textbf{Розділ SRS} & \textbf{Критерій прийняття} & \textbf{BDD-сценарій} \\
\midrule
\endhead
\bottomrule
\endfoot
\bottomrule
\endlastfoot
REQ-F-001 & K2, K22 & \S\ref{sec:funktsional-katalog} & Публічний перегляд фільмів і сеансів без автентифікації & --- \\
REQ-F-002 & K2, K31 & \S\ref{sec:funktsional-katalog} & Відображення схеми залу; вибір місця не резервує його & --- \\
REQ-F-003 & K23 & \S\ref{sec:funktsional-auth} & Бронювання та скасування суто після автентифікації & --- \\
REQ-F-004 & K24 & \S\ref{sec:funktsional-auth} & Унікальний числовий внутрішній ID користувача & --- \\
REQ-F-005 & K24 & \S\ref{sec:funktsional-auth} & Доступ виключно до власних бронювань (захист IDOR) & --- \\
REQ-F-006 & K26 & \S\ref{sec:funktsional-auth} & Відокремлений адміністративний розділ за моделлю RBAC & --- \\
REQ-F-007 & K4 & \S\ref{sec:funktsional-bronuvannia} & Обмеження кількості місць: від 1 до 6 включно & --- \\
REQ-F-008 & K5 & \S\ref{sec:funktsional-bronuvannia} & Атомарне створення в момент підтвердження без hold & Сценарій 1 \\
REQ-F-009 & K29 & \S\ref{sec:funktsional-bronuvannia} & Рівно одне успішне бронювання при конкуренції за місце & Сценарій 1 \\
REQ-F-010 & K30 & \S\ref{sec:funktsional-bronuvannia} & Відхилення бронювання, якщо хоч одне місце зайняте & Сценарій 2 \\
REQ-F-011 & K36 & \S\ref{sec:funktsional-skasuvannia-user} & Скасування лише всього бронювання повністю (без часткового) & Сценарій 3 \\
REQ-F-012 & K37 & \S\ref{sec:funktsional-skasuvannia-user} & Скасування дозволене виключно до настання часу сеансу & Сценарії 3, 4, 5 \\
REQ-F-013 & K38 & \S\ref{sec:funktsional-skasuvannia-user} & Перехід у статус «Скасоване» та негайне звільнення місць & Сценарій 3 \\
REQ-F-014 & K39 & \S\ref{sec:funktsional-skasuvannia-user} & Ідемпотентність повторних спроб скасування & --- \\
REQ-F-015 & K3 & \S\ref{sec:funktsional-admin} & Управління довідниками фільмів, залів і сеансів & --- \\
REQ-F-016 & K44 & \S\ref{sec:funktsional-admin} & Вільне редагування параметрів сеансу без бронювань & --- \\
REQ-F-017 & K45 & \S\ref{sec:funktsional-admin} & Блокування прямого редагування сеансу з бронюваннями & Сценарій 6 \\
REQ-F-018 & K46 & \S\ref{sec:funktsional-admin} & Явне скасування сеансу адміном з обов'язковою причиною & Сценарій 6 \\
REQ-F-019 & K47 & \S\ref{sec:funktsional-admin} & Каскадний перехід бронювань у «Скасовано» та збереження & Сценарій 6 \\
REQ-F-020 & K48 & \S\ref{sec:funktsional-admin} & Відображення причини скасування сеансу в кабінеті & Сценарій 6 \\
REQ-TDD-001 & K7, K11 & \S\ref{sec:tdd-process} & Обов'язковість попереднього створення автотестів & --- \\
REQ-TDD-002 & K12 & \S\ref{sec:tdd-process} & Дотримання 5-крокового циклу Red-Green-Refactor & --- \\
REQ-TDD-003 & K13 & \S\ref{sec:tdd-scope} & Покриття тестами всіх погоджених граничних кейсів & Сценарії 1--6 \\
REQ-TDD-004 & K14 & \S\ref{sec:tdd-scope} & Якісна перевірка правил без штучного порогу покриття & --- \\
REQ-TDD-005 & K15, K43 & \S\ref{sec:tdd-scope} & Вимоги до тестового раннера (швидкість, емуляція) & --- \\
REQ-LOG-001 & K8, K16 & \S\ref{sec:loguvannia-podii} & Автоматичне логування 5 обов'язкових категорій подій & --- \\
REQ-LOG-002 & K17 & \S\ref{sec:loguvannia-podii} & Наявність обов'язкових полів аудиту в кожному записі & --- \\
REQ-LOG-003 & K18 & \S\ref{sec:loguvannia-bezpeka} & Сувора заборона збереження паролів і токенів у логах & --- \\
REQ-LOG-004 & K19 & \S\ref{sec:loguvannia-zberezhennia} & Персистентність та збереження логів після рестарту & --- \\
REQ-LOG-005 & K20 & \S\ref{sec:loguvannia-testy} & Збереження окремих звітів кожного прогону тестів TDD & --- \\
REQ-NF-001 & K5, K30 & \S\ref{sec:nefunktsional-atomarnist} & Транзакційна атомарність бронювання (ACID) & Сценарії 1, 2 \\
REQ-NF-002 & K40, K41 & \S\ref{sec:nefunktsional-atomarnist} & Персистентність та збереження узгодженості даних & --- \\
REQ-NF-003 & K40 & \S\ref{sec:nefunktsional-bezpeka} & Безумовний серверний контроль прав доступу RBAC & --- \\
REQ-NF-004 & K24 & \S\ref{sec:nefunktsional-bezpeka} & Ізоляція даних бронювань між користувачами (IDOR) & --- \\
REQ-NF-005 & K27 & \S\ref{sec:nefunktsional-bezpeka} & Заборона самопризначення ролі адміністратора клієнтом & --- \\
REQ-NF-006 & K29, K43 & \S\ref{sec:nefunktsional-konkurentnist} & Коректна конкурентна обробка без подвійних бронювань & Сценарій 1 \\
\bottomrule
\end{longtable}
\end{small}


\end{document}
